\documentclass[../main.tex]{subfiles}

\begin{document}

%% ======================================================================
%%  Section 4 – HDF5 Extensions
%% ======================================================================
\section{HDF5 Library Extensions}
\label{sec:extensions}

\subsection{Assessment Boundary}

Supply-chain issues are treated in Section~\ref{sec:supply-chain}. This
chapter addresses a different question: how VOL connectors, filter plugins,
and Virtual File Drivers (VFDs) can extend a deployment's privacy boundary
beyond the logical HDF5 file.

The P1--P7 labels use the HDF5 Privacy Exposure Model summarized in the core
chapter: P1 semantic metadata, P2 structural or on-disk leakage, P3 raw-data
patterns, P4 unintended inclusion during processing, P5 aggregation or
inference, P6 persistence beyond intended lifetime, and P7 misplaced trust in
privacy by format. \code{EXT}, \code{OPS}, \code{TCD}, and \code{SCD} are
separate SSP tags for extensions, operations, toolchains, and distribution.
These labels classify candidate mechanisms, not inherent risk or response
urgency. Without a bounded deployment, data population, recipient, retention
condition, and consequence, the chapter-wide observation remains unrated
evidence-backlog item \code{EXT-PRIV-001}.

\begin{table}[htbp]
\centering
\small
\setlength{\tabcolsep}{3pt}
\renewcommand{\arraystretch}{1.12}
\begin{tabularx}{\textwidth}{@{}p{0.17\textwidth}p{0.25\textwidth}X p{0.18\textwidth}@{}}
\hline
\textbf{Extension} & \textbf{Position in workflow} &
\textbf{Boundary added to the assessment} & \textbf{Candidate families} \\
\hline
VOL connector & Above or in place of native object storage &
External and non-HDF5 sources, aggregation, generated objects and metadata,
network services, credentials, caches, and connector code. &
P1, P4--P7; SSP: EXT, OPS, SCD \\
Filter & Immediately before storage and after retrieval &
Pipeline metadata, compressed-size patterns, plaintext buffers, derived
information, temporary artifacts, and dynamically loaded code. &
P1--P7; SSP: EXT, SCD \\
VFD & Physical byte storage and retrieval &
Storage locations, replicas and histories, physical layout and access
patterns, network and cloud infrastructure, and VFD code. &
P2--P7; SSP: EXT, OPS, SCD \\
\hline
\end{tabularx}
\caption{Extension mechanisms and the deployment boundary they can add}
\label{tab:extension-boundary}
\end{table}

\subsection{Virtual Object Layer Connectors}

A connector, or a connector stack, can present external and heterogeneous
resources through the HDF5 API. Protecting or inspecting the apparent file is
therefore not sufficient to establish where all data originated, where they
were processed or retained, or which policies applied. Table~\ref{tab:vol-privacy-risks}
consolidates the mechanisms that must be tested in a concrete deployment.

\begingroup
\small
\setlength{\tabcolsep}{3pt}
\renewcommand{\arraystretch}{1.14}
\begin{longtable}{@{}p{0.20\textwidth}p{0.62\textwidth}p{0.13\textwidth}@{}}
\caption{Candidate privacy-exposure families for VOL connector mechanisms}
\label{tab:vol-privacy-risks}\\
\hline
\textbf{Mechanism} & \textbf{Deployment-relevant observation} &
\textbf{Families / tags} \\
\hline
\endfirsthead
\caption[]{Candidate privacy-exposure families for VOL connector mechanisms, continued}\\
\hline
\textbf{Mechanism} & \textbf{Deployment-relevant observation} &
\textbf{Families / tags} \\
\hline
\endhead
\hline
\endfoot

External and non-HDF5 sources &
A connector may retrieve and integrate cloud objects, databases, remote REST
or other network services, live streams, TIFF/GeoTIFF, BUFR, GRIB2, CDF,
DAOS, or other sources, and may generate derived data or virtual objects.
Their origin, location, ownership, and governing policy need not be visible to
an application that sees one logical HDF5 dataset. &
P7; SSP: EXT \\
\hline

Aggregation and policy crossing &
Independent sources can be combined without an explicit application-level
policy decision. For example, geospatial records, near-real-time weather and
disaster feeds, and property assessments could appear as one property-impact
dataset even though ownership, access, retention, or privacy rules differ. &
P5, P7; SSP: EXT, OPS \\
\hline

Generated metadata and provenance &
Relationships, locations, URLs, bucket names, origin, timestamps,
organizational structure, and modification history can disclose context while
raw data remain protected. BUFR and GRIB2 connectors may create file
inventories, derived attributes, and dimension scales. &
P1, P4; SSP: EXT \\
\hline

Network and cache artifacts &
Headers, DNS and proxy logs, client history, authentication records, local
caches, temporary downloads, raw data, metadata, and credentials can record
which datasets were used and may remain after execution or be copied and
archived. Networked connectors therefore inherit the surrounding network's
privacy and security properties. &
P4, P6; SSP: EXT, OPS \\
\hline

In-process plugin trust &
A dynamically loaded connector executes in the application's process and can
inspect, modify, copy, cache, export, or share data before returning. Its code,
provenance, search path, and configuration belong in the trusted computing
base. &
P4, P6; SSP: EXT, SCD \\
\hline
\end{longtable}
\endgroup

\subsection{Filter Plugins}

Filters process data immediately before storage and immediately after
retrieval. Their behavior, execution environment, pipeline metadata, and
artifacts are therefore part of the privacy assessment. Table~\ref{tab:filter-plugin-privacy}
combines the cross-cutting mechanisms with representative built-in and
third-party filters; it does not assign risk to a filter merely because a
mechanism is possible.

\begingroup
\small
\setlength{\tabcolsep}{3pt}
\renewcommand{\arraystretch}{1.12}
\begin{longtable}{@{}p{0.23\textwidth}p{0.59\textwidth}p{0.13\textwidth}@{}}
\caption{Primary privacy mechanisms associated with HDF5 filters}
\label{tab:filter-plugin-privacy}\\
\hline
\textbf{Mechanism or filters} & \textbf{Representative consideration} &
\textbf{Families / tags} \\
\hline
\endfirsthead
\caption[]{Primary privacy mechanisms associated with HDF5 filters, continued}\\
\hline
\textbf{Mechanism or filters} & \textbf{Representative consideration} &
\textbf{Families / tags} \\
\hline
\endhead
\hline
\endfoot

Pipeline identifiers and parameters &
The object header records the pipeline required to interpret stored data.
Identifiers and settings may reveal a security mechanism, organizational
standard, software ecosystem, application domain, precision, resolution,
floating-point accuracy, loss tolerance, or experimental assumptions. This
metadata generally cannot be removed without affecting interoperability. &
P1, P3; SSP: EXT \\
\hline

Deflate (GZIP), Zstandard, LZ4, and BZIP2 &
Method, level, block choice, and resulting chunk sizes can reveal entropy,
sparsity, redundancy, and structure. LZ4's fast, modest-ratio behavior and
BZIP2 block characteristics are instances of the same observable. &
P1--P3; SSP: EXT \\
\hline

Shuffle and Blosc &
Reordering is reversible, but its use with compression, block sizes,
shuffle, or bitshuffle options can disclose numeric representation and data
organization. &
P1, P3; SSP: EXT \\
\hline

Fletcher32 &
Checksum metadata reveals that integrity protection is used but generally
introduces little additional privacy exposure. &
P1; SSP: EXT \\
\hline

N-bit and Scale-Offset &
Bit width, scale, and offset can reveal effective precision, expected value
ranges, resolution, and acceptable quantization error. &
P1, P3; SSP: EXT \\
\hline

ZFP, SZ, and SZ3 &
Fixed-rate, fixed-precision, fixed-accuracy modes and error bounds disclose
acceptable information loss and numerical-accuracy requirements. &
P1, P3; SSP: EXT \\
\hline

Plaintext processing and derived information &
Filters see plaintext buffers before protective compression or encryption and
may compute or retain statistics, patterns, sparsity, or other derived
information. Buffers may remain in memory, swap, or crash dumps; implementations
may also create logs or temporary files. &
P4--P7; SSP: EXT \\
\hline

Dynamic loading and plugin trust &
A filter loaded from an environment-controlled search directory executes in
the application process and may inspect memory, copy decoded data, log,
retain, modify, or transmit it. Search results can differ between systems;
untrusted or user-writable paths therefore create a supply-chain boundary. &
P4, P6, P7; SSP: EXT, SCD \\
\hline
\end{longtable}
\endgroup

\subsection{Virtual File Drivers}

VFDs do not change the logical HDF5 data model, but determine where and how
bytes are stored and retrieved. A logical file may reside in cloud object
storage, a distributed file system, network storage, process memory, or across
multiple devices, so its apparent name need not identify the governing storage
or policy domain. Consequently, physical layout, copies, communications,
storage-provider controls, and artifacts can matter even when logical HDF5
objects are adequately protected.

\begingroup
\small
\setlength{\tabcolsep}{3pt}
\renewcommand{\arraystretch}{1.12}
\begin{longtable}{@{}p{0.20\textwidth}p{0.62\textwidth}p{0.13\textwidth}@{}}
\caption{Candidate privacy-exposure families for HDF5 Virtual File Drivers}
\label{tab:vfd-privacy-risks}\\
\hline
\textbf{VFD or mechanism} & \textbf{Representative consideration} &
\textbf{Families / tags} \\
\hline
\endfirsthead
\caption[]{Candidate privacy-exposure families for HDF5 Virtual File Drivers, continued}\\
\hline
\textbf{VFD or mechanism} & \textbf{Representative consideration} &
\textbf{Families / tags} \\
\hline
\endhead
\hline
\endfoot

SEC2 and STDIO &
Persistent storage can expose residual information in reused blocks if the
storage layer does not securely initialize or overwrite them. &
P6 \\
\hline

Core &
The file resides in process memory and may persist in swap or crash dumps; with
the backing-store option, it is written to persistent storage at close. &
P6 \\
\hline

Family, Split, Multi, and Subfiling &
One logical file becomes multiple independently managed files. Family and
Subfiling divide data across members or devices; Split separates metadata and
raw data; Multi separates HDF5 memory types. Different permissions, locations,
or incomplete deletion can disclose fragments, metadata, or raw data. &
P2, P6, P7; SSP: EXT \\
\hline

Mirror, MPI-I/O, and ROS3 &
Mirror replicates to a remote service and depends on communication,
authentication, daemon trust, and protection of copies. MPI-I/O passes data
through HPC communications unless protected below HDF5. ROS3 uses
S3-compatible stores; HTTP requests and Range requests, DNS, and cloud access
logs create artifacts and disclose access. &
P4, P6, P7; SSP: EXT, OPS, SCD \\
\hline

Log (VFD SWMR) and Onion &
Metadata journals retain snapshots for recovery. Onion retains historical
versions, allowing recovery of deleted or overwritten data and reconstruction
of metadata evolution, provenance, and workflow history. &
P2, P5, P6; SSP: EXT \\
\hline

Physical layout and access patterns &
Offsets, allocations, file organization, access frequency, read/write ratios,
sequential or random I/O, update frequency, and touched regions can reveal
dataset, algorithm, storage-architecture, or user characteristics to operating
systems, distributed stores, servers, or cloud providers. &
P2, P3; SSP: EXT \\
\hline

Workflow artifacts and plugin trust &
Temporary files, swap, memory images, crash dumps, replicas, journals, and
cloud logs can persist outside the logical file. A dynamically loaded VFD can
inspect, modify, retain, duplicate, or export data in process and therefore
belongs in the trusted computing base. &
P4, P6; SSP: EXT, SCD \\
\hline
\end{longtable}
\endgroup

\subsection{Control Proposals and Routing}

The mechanisms above support control design but do not silently promote
\code{EXT-PRIV-001} into a formal finding. The authoritative recommendation
register and phased roadmap provide ownership, dependencies, sequencing, and
acceptance evidence. The identifiers below express neither rank nor urgency.

\begingroup
\small
\setlength{\tabcolsep}{3pt}
\renewcommand{\arraystretch}{1.11}
\begin{longtable}{@{}p{0.14\textwidth}p{0.81\textwidth}@{}}
\caption{Extension privacy and trust recommendations}\\
\hline
\textbf{Identifier} & \textbf{Control outcome} \\
\hline
\endfirsthead
\caption[]{Extension privacy and trust recommendations, continued}\\
\hline
\textbf{Identifier} & \textbf{Control outcome} \\
\hline
\endhead
\hline
\endfoot

\code{VOL-S1} & Document external sources, generated metadata, caching,
network communications, authentication, workflow artifacts, and privacy
assumptions for each connector. \\
\hline
\code{VOL-S2} & Standardize introspection for connector name, version, stack,
source types and locations, network use, cached data or credentials, temporary
or persistent artifacts, generated or derived data, cross-policy aggregation,
and privacy configuration. Existing identity and capability callbacks do not
expose all of these characteristics. \\
\hline
\code{VOL-S3} & Treat connectors as trusted code: restrict search paths,
exclude user-writable locations, and verify identity, provenance, and
integrity; support enforceable loading policy. \\
\hline
\code{VOL-S4} & Minimize unnecessary generated metadata and persistent caches,
unless explicitly requested, and document unavoidable privacy characteristics
during connector design. \\
\hline

\code{FILTER-S1} & Apply the same trusted-plugin, restricted-search-path,
provenance, integrity, and enforceable-loading controls to filters. \\
\hline
\code{FILTER-S2} & Select and document pipelines with explicit consideration
of what identifiers, parameters, precision, compression, and workflow choices
reveal. \\
\hline
\code{FILTER-S3} & Where a file recipient must not see pipeline or structural
metadata, evaluate whole-file protection covering raw data and metadata. Key
management and recovery evidence follows \code{RM-2G}. Such protection can
prevent that recipient from inspecting filter identifiers, parameters, object
headers, and structural metadata, but does not treat caches, logs, authorized
misuse, post-decryption inference, retention, or key-loss risk. \\
\hline
\code{FILTER-S4} & Minimize workflow artifacts, erase temporary buffers when
possible using an appropriate secure method, avoid persistent temporary files,
and document remaining artifacts. \\
\hline
\code{FILTER-S5} & Standardize filter documentation for generated metadata,
plaintext processing, derived information, temporary artifacts, and privacy
implications of parameters. \\
\hline

\code{VFD-S1} & Document storage locations, additional copies, network use,
workflow artifacts, and storage assumptions for every VFD. \\
\hline
\code{VFD-S2} & Give replicated files, journals, histories, and backing stores
secure deletion, retention, and lifecycle controls. \\
\hline
\code{VFD-S3} & Document authentication, communications, logging, and
storage-provider privacy assumptions for cloud and distributed VFDs. \\
\hline
\code{VFD-S4} & Apply trusted-plugin, restricted-search-path, provenance,
integrity, and enforceable-loading controls to dynamically loaded VFDs. \\
\hline
\end{longtable}
\endgroup

The trusted-loading proposals call for configurable policy based on plugin
identity, exposed characteristics, or administrative rules, in addition to
deployment-side search-path and provenance controls.

For \code{VOL-S3}, \code{FILTER-S1}, and \code{VFD-S4}, signed-plugin
verification requires HDF5 2.2.0 or a later verified release built with
\code{HDF5_REQUIRE_SIGNED_PLUGINS=ON}, a configured and protected trusted
keystore, signed approved artifacts, and a test that unsigned or untrusted
plugins are rejected. Upgrading alone does not enable the control
(Table~\ref{tab:tools-tests-security-features}).

No report-wide risk urgency is assigned to \code{EXT-PRIV-001}. Phase~2
\code{RM-2E} establishes a bounded deployment, data flow, recipient, retention
condition, and consequence; that administrative target is not a risk-response
clock or a Low-risk conclusion. \code{RM-3D} and \code{RM-3F} then provide
proactive extension assurance and evidence-dependent privacy treatment. A
deployment matching separately rated plugin-search finding
\code{APP-SEC-005} instead follows that finding's Near-term application clock;
it does not assign a High tier to extensions as a class.

\end{document}
